\section{Síntesis y ampliación}

\subsection{Conclusiones centrales}

\begin{enumerate}
\item Lo central de la socialización con IA no es que la IA converse con el usuario para hacerle compañía, sino que el doble digital, con su context, cree por iniciativa propia conexiones entre personas reales.
\item El context es el activo central de la socialización con IA: la inteligencia quizá llegue a estar al alcance de todos por igual; el context, no.
\item La subjetividad es la expresión en ingeniería de «ser uno mismo» y es la base para que el doble digital pueda actuar en representación del usuario.
\item La privacidad y el context forman una relación de intercambio: el producto debe hacer que el usuario perciba un beneficio antes de abandonarlo.
\item Elys no es un Tinder con IA; su objetivo es el emparejamiento por context de alta dimensión y la socialización proactive.
\item Las empresas de modelos no ganan por naturaleza la socialización con IA; la ventaja defensiva de las empresas de aplicaciones está en el context, las relaciones y el lugar que el producto ocupa en la mente del usuario.
\end{enumerate}

\subsection{Preguntas abiertas}

\begin{itemize}
\item ¿Estarán dispuestos los usuarios a entregar suficiente context a una empresa comercial?
\item Cuando el doble digital se expresa en nombre del usuario, ¿cómo se establecen los mecanismos de autorización, revocación y responsabilidad?
\item ¿La socialización con IA aumentará las conexiones entre personas reales o producirá más interacciones ficticias?
\item Entre actores como ChatGPT, Meta, Doubao (豆包) y Elys, ¿quién tiene más probabilidades de obtener una ventaja de context a largo plazo?
\item ¿Cómo puede un producto proactive evitar convertirse en acoso o salirse de control?
\end{itemize}

\subsection{Lecturas adicionales}

\begin{itemize}
\item EP139, Historia técnica de los Agent: para entender el contexto técnico de los proactive agent.
\item EP136, El modelo como OS: para entender cómo el modelo se convierte en la capa de planificación de tareas.
\item Material sobre interacción humano-computadora, redes sociales, sistemas de recomendación y cómputo con preservación de la privacidad: para entender las bases multidisciplinarias de las redes sociales con IA.
\item La película \textit{Her} y el debate sobre los productos de AI companion: para entender la diferencia entre la narrativa de la compañía y los productos reales.
\end{itemize}

\begin{importantbox}{Valoración final}
La verdadera dificultad de la socialización con IA no es lograr que la IA sepa conversar, sino lograr que, respetando la privacidad y la subjetividad, convierta el context de una persona en mejores conexiones con personas reales. Si lo consigue, la unidad básica de las redes sociales pasará de ser la «cuenta» a ser el «doble digital capaz de actuar».\end{importantbox}

\end{document}
